\chapter{Bibliographie}\label{ch:bibliographie}

La bibliographie réunit les sources numériques, les recueils, les instruments et les études cités dans le volume. Les références d’édition propres à chaque inscription sont données dans le catalogue et dans les concordances. Pour chaque titre, une mention indique le niveau de contrôle de la référence : « vérifié » signifie que la notice bibliographique (auteur, titre, lieu, date, collection, pagination ou DOI lorsqu’ils sont donnés) a été contrôlée sur un catalogue de bibliothèque, sur le site de l’éditeur ou sur une base de métadonnées ; « contenu consulté » signifie que le texte lui-même a été lu pour ce volume. Aucune vérification de notice ne vaut lecture de l’ouvrage, et les passages des études ne sont jamais invoqués sans que le texte ait été consulté.

\newcommand{\bib}[2]{\par\noindent\hangindent=1.5em #1\par\noindent\hangindent=1.5em\hspace{1.5em}{\footnotesize\color{prismegris}#2}\par\smallskip}

\section*{Sources et bases de données}\addcontentsline{toc}{section}{Sources et bases de données}

\bib{\emph{Epigraphik-Datenbank Clauss/Slaby} (EDCS), M.~Clauss, A.~Kolb, W.~A.~Slaby, B.~Woitas : db.edcs.eu.}{Consultée à travers les exports de 2020 et de 2022 publiés par le projet SDAM ; texte de référence du catalogue (« texte transmis »).}

\bib{SDAM (Social Dynamics and Complexity in the Ancient Mediterranean, Aarhus University), \emph{EDCS\_ETL}, dépôt de code et de données, fichiers \emph{2020-12-09-province\_Roma} et \emph{2022-09-12-EDCS\_via\_Lat\_Epig-prov\_Roma} : github.com/sdam-au/EDCS\_ETL.}{Contenu consulté : clonage du dépôt, octobre 2026.}

\bib{C.~Carletti, A.~E.~Felle, A.~Rocco, \emph{EDB. Epigraphic Database Bari — Inscriptions by Christians in Rome (3rd–8th cent. CE). EpiDoc files}, Zenodo, 2019, DOI 10.5281/zenodo.3575485.}{Contenu consulté : archive EpiDoc téléchargée et analysée intégralement (441 fiches appariées au catalogue). Licence du dépôt : CC BY 4.0 ; licence des fichiers : CC BY-NC-SA 3.0.}

\bib{\emph{Epigraphic Database Bari} (EDB) : www.edb.uniba.it.}{Les textes des listes de résultats ont été utilisés à travers le fichier de collation du projet ; le site n’a pas été interrogé directement (accès automatisé refusé).}

\bib{\emph{Epigraphic Database Heidelberg} (EDH) : edh.ub.uni-heidelberg.de.}{Fiches HD009698, HD021839, HD028015 et HD030246 consultées pour les doublons et le \nnum{98}.}

\bib{Center for Epigraphical and Palaeographical Studies, The Ohio State University, collection photographique des ILCV (estampage ILCV 3252A ; photographie du \nnum{68}).}{Contenu consulté : deux images examinées ; aucune reproduction publiée dans ce volume.}

\bib{\emph{The Cult of Saints in Late Antiquity}, University of Oxford, base de données en ligne, notice E07620.}{Information de seconde main, citée au chapitre~\ref{ch:historiographie} ; notice non réexaminée.}

\bib{\emph{EpiDoc Guidelines: Ancient Documents in TEI XML} : epidoc.stoa.org/gl/.}{Documentation consultée pour les conventions éditoriales et la correspondance avec le système Krummrey–Panciera.}

\section*{Recueils épigraphiques}\addcontentsline{toc}{section}{Recueils épigraphiques}

\bib{G.~B.~de Rossi, \emph{Inscriptiones christianae urbis Romae septimo saeculo antiquiores}, I, Rome, 1857–1861 ; II, 1, Rome, 1888 ; supplément, éd. G.~Gatti, Rome, 1915.}{Vérifié (HathiTrust, Internet Archive) ; supplément : notices de catalogue.}

\bib{\emph{Inscriptiones christianae urbis Romae septimo saeculo antiquiores}, nova series, I–X, Rome / Cité du Vatican, 1922–1992 : I (A.~Silvagni, 1922) ; II (A.~Silvagni, 1935) ; III (A.~Silvagni, A.~Ferrua, 1956) ; IV–VIII (A.~Ferrua, 1964–1983) ; IX (A.~Ferrua, D.~Mazzoleni, 1985) ; X (D.~Mazzoleni, C.~Carletti, 1992).}{Vérifié (catalogues de La Sapienza, de Stanford et de RERO ; liste des publications du PIAC). Les volumes n’ont pas été consultés : les lectures du catalogue sont celles des bases.}

\bib{E.~Diehl, \emph{Inscriptiones Latinae Christianae Veteres}, I–III, Berlin, Weidmann, 1925–1931 ; IV, \emph{Supplementum}, éd. J.~Moreau, H.-I.~Marrou, Dublin – Zurich, Weidmann, 1967.}{Vérifié (HathiTrust, Stanford SearchWorks). Les notes de Diehl citées le sont à travers l’EDB.}

\bib{A.~Ferrua, \emph{Epigrammata Damasiana}, Cité du Vatican, Pontificio Istituto di Archeologia Cristiana, 1942 (Sussidi allo studio delle antichità cristiane, 2).}{Vérifié (liste des publications du PIAC).}

\bib{F.~Bücheler, E.~Lommatzsch, \emph{Carmina Latina Epigraphica}, Leipzig, 1895–1926.}{Cité d’après les concordances des bases ; notice non contrôlée.}

\bib{\emph{Corpus Inscriptionum Latinarum}, VI et XIV.}{Cité d’après les concordances des bases.}

\bib{\emph{L’Année épigraphique}, Paris, 1888–.}{Cité d’après les concordances des bases.}

\section*{Instruments}\addcontentsline{toc}{section}{Instruments}

\bib{R.~S.~Bagnall, A.~Cameron, S.~R.~Schwartz, K.~A.~Worp, \emph{Consuls of the Later Roman Empire}, Atlanta, Scholars Press, 1987 (Philological Monographs of the American Philological Association, 36), viii–759 p.}{Vérifié (Oxford University Press, FRANTIQ, HathiTrust) ; contenu non consulté. Les renvois à la CLRE sont faits à travers l’EDB.}

\bib{A.~H.~M.~Jones, J.~R.~Martindale, J.~Morris, \emph{The Prosopography of the Later Roman Empire}, I, \emph{A.D.~260–395}, Cambridge, 1971.}{Notice vérifiée ; contenu non consulté.}

\bib{Ch.~Pietri, L.~Pietri (dir.), \emph{Prosopographie chrétienne du Bas-Empire}, 2, \emph{Prosopographie de l’Italie chrétienne (313–604)}, Rome, École française de Rome, 1999–2000, 2 vol.}{Vérifié (Persée ; notices de libraires) ; contenu non consulté.}

\bib{V.~Grumel, \emph{La chronologie}, Paris, 1958 (Traité d’études byzantines, I).}{Notice vérifiée ; contenu non consulté.}

\bib{H.~Krummrey, S.~Panciera, « Criteri di edizione e segni diacritici », \emph{Tituli}, 2, 1980, p.~205–215.}{Référence et valeur des signes vérifiées sur les EpiDoc Guidelines.}

\bib{B.~Rhodes, \emph{PyEphem}, version 4.2.1, bibliothèque de calcul astronomique.}{Logiciel utilisé pour les calculs lunaires (chapitre~\ref{ch:calendrier}).}

\section*{Études}\addcontentsline{toc}{section}{Études}

\bib{J.~N.~Adams, \emph{Social Variation and the Latin Language}, Cambridge, Cambridge University Press, 2013.}{Notice vérifiée ; contenu non consulté.}

\bib{J.~Bodel (éd.), \emph{Epigraphic Evidence. Ancient History from Inscriptions}, Londres – New York, 2001.}{Notice vérifiée ; contenu non consulté.}

\bib{P.~Brown, \emph{The Cult of the Saints. Its Rise and Function in Latin Christianity}, Chicago, University of Chicago Press, 1981 (Haskell Lectures on History of Religions, n.s., 2).}{Partiellement vérifié (Internet Archive).}

\bib{C.~Bruun, J.~Edmondson (éd.), \emph{The Oxford Handbook of Roman Epigraphy}, Oxford – New York, 2015.}{Notice vérifiée ; contenu non consulté.}

\bib{C.~Carletti, \emph{Iscrizioni cristiane di Roma. Testimonianze di vita cristiana (secoli III–VII)}, Florence, 1986.}{Notice vérifiée ; contenu non consulté.}

\bib{C.~Carletti, \emph{Epigrafia dei cristiani in Occidente dal III al VII secolo. Ideologia e prassi}, Bari, Edipuglia, 2008.}{Vérifié ; contenu non consulté.}

\bib{M.~Carroll, \emph{Spirits of the Dead. Roman Funerary Commemoration in Western Europe}, Oxford, 2006.}{Notice vérifiée ; contenu non consulté.}

\bib{S.~Diefenbach, \emph{Römische Erinnerungsräume. Heiligenmemoria und kollektive Identitäten im Rom des 3. bis 5. Jahrhunderts n.~Chr.}, Berlin – New York, De Gruyter, 2007 (Millennium-Studien, 11), DOI 10.1515/9783110204261.}{Vérifié (De Gruyter) ; contenu non consulté.}

\bib{R.~Duncan-Jones, \emph{Structure and Scale in the Roman Economy}, Cambridge, 1990.}{Notice vérifiée ; contenu non consulté.}

\bib{A.~E.~Felle, « Inscriptions by Christians in Late Antique Rome: Some Issues and Perspectives for the Epigraphic Database Bari (EDB) », dans I.~Velázquez Soriano, D.~Espinosa Espinosa (éd.), \emph{Epigraphy in the Digital Age}, Oxford, Archaeopress, 2021, p.~213–224.}{Notice vérifiée (catalogue de la TU München) ; contenu non consulté.}

\bib{V.~Fiocchi Nicolai, F.~Bisconti, D.~Mazzoleni, \emph{Le catacombe cristiane di Roma. Origini, sviluppo, apparati decorativi, documentazione epigrafica}, Ratisbonne, Schnell \& Steiner, 1998.}{Notice vérifiée ; éditeur confirmé par la 4\textsuperscript{e} édition (2025, ISBN 978-3-7954-3967-5) ; contenu non consulté.}

\bib{F.~Grossi Gondi, \emph{Trattato di epigrafia cristiana latina e greca del mondo romano occidentale}, Rome, 1920.}{Notice vérifiée ; contenu non consulté.}

\bib{M.~A.~Handley, \emph{Death, Society and Culture. Inscriptions and Epitaphs in Gaul and Spain, AD~300–750}, Oxford, 2003 (BAR International Series, 1135).}{Notice vérifiée ; contenu non consulté.}

\bib{P.~Heřmánková, V.~Kaše, A.~Sobotková, « Inscriptions as Data: Digital Epigraphy in Macro-historical Perspective », \emph{Journal of Digital History}, 1/1, 2021, DOI 10.1515/jdh-2021-1004.}{Texte consulté ; DOI vérifié (Crossref).}

\bib{K.~Hopkins, « On the Probable Age Structure of the Roman Population », \emph{Population Studies}, 20/2, 1966, p.~245–264, DOI 10.1080/00324728.1966.10406097.}{Vérifié (PubMed, PMID 22084913) ; contenu non consulté.}

\bib{M.~Humphries, « The Memory of Mursa: Usurpation, Civil War, and Contested Legitimacy Under the Sons of Constantine », dans N.~Baker-Brian, S.~Tougher (éd.), \emph{The Sons of Constantine, AD~337–361}, Cham, Palgrave Macmillan, 2020, p.~157–183, DOI 10.1007/978-3-030-39898-9\_6.}{Métadonnées vérifiées ; contenu non consulté.}

\bib{I.~Kajanto, \emph{Onomastic Studies in the Early Christian Inscriptions of Rome and Carthage}, Helsinki, 1963 (Acta Instituti Romani Finlandiae, II, 1).}{Notice vérifiée ; contenu non consulté.}

\bib{I.~Kajanto, \emph{The Latin Cognomina}, Helsinki, 1965.}{Notice vérifiée ; contenu non consulté.}

\bib{R.~MacMullen, « The Epigraphic Habit in the Roman Empire », \emph{American Journal of Philology}, 103/3, 1982, p.~233–246, DOI 10.2307/294470.}{Métadonnées vérifiées ; contenu non consulté.}

\bib{D.~Mazzoleni, \emph{Epigrafi del mondo cristiano antico}, Rome, 2002.}{Notice vérifiée ; contenu non consulté.}

\bib{E.~A.~Meyer, « Explaining the Epigraphic Habit in the Roman Empire: The Evidence of Epitaphs », \emph{Journal of Roman Studies}, 80, 1990, p.~74–96, DOI 10.2307/300281.}{Vérifié (Cambridge Core) ; contenu non consulté.}

\bib{H.~Nordberg, \emph{Biometrical Notes. The Information on Ancient Christian Inscriptions from Rome concerning the Duration of Life and the Dates of Birth and Death}, Helsinki, 1963 (Acta Instituti Romani Finlandiae, II, 2).}{Notice vérifiée ; contenu non consulté.}

\bib{J.~J.~O’Donnell, « The Career of Virius Nicomachus Flavianus », \emph{Phoenix}, 32/2, 1978, p.~129–143, DOI 10.2307/1087470.}{Métadonnées vérifiées ; contenu non consulté.}

\bib{Ch.~Pietri, \emph{Roma christiana. Recherches sur l’Église de Rome, son organisation, sa politique, son idéologie de Miltiade à Sixte III (311–440)}, Rome, École française de Rome, 1976 (BEFAR, 224), 2 vol., DOI 10.3406/befar.1976.1278.}{Vérifié (Persée, WorldCat) ; pagination totale divergente selon les sources, non donnée ; contenu non consulté.}

\bib{Ch.~Pietri, « Appendice prosopographique à la \emph{Roma christiana} (311–440) », \emph{Mélanges de l’École française de Rome. Antiquité}, 89/1, 1977, p.~371–415.}{Vérifié (Persée) ; contenu non consulté.}

\bib{Ch.~Pietri, « Remarques sur l’onomastique chrétienne de Rome », dans \emph{L’onomastique latine}, Paris, CNRS, 1977, p.~437–456 ; repris dans \emph{Christiana respublica}, 1997, p.~1571–1582.}{Partiellement vérifié (Persée ; une autre source donne p.~437–445) ; contenu non consulté.}

\bib{Ch.~Pietri, « Grabinschrift II (lateinisch) », \emph{Reallexikon für Antike und Christentum}, 12, Stuttgart, 1983, col.~514–590 ; version française dans \emph{Christiana respublica}, 1997, p.~1407–1468.}{Vérifié ; contenu non consulté.}

\bib{Ch.~Pietri, \emph{Christiana respublica. Éléments d’une enquête sur le christianisme antique}, Rome, École française de Rome, 1997 (Collection de l’École française de Rome, 234), 3 vol.}{Vérifié (BnF, Persée) ; contenu non consulté.}

\bib{É.~Rebillard, \emph{Religion et sépulture. L’Église, les vivants et les morts dans l’Antiquité tardive}, Paris, Éditions de l’EHESS, 2003 (Civilisations et sociétés, 115) ; trad. angl. \emph{The Care of the Dead in Late Antiquity}, Ithaca, Cornell University Press, 2009 (Cornell Studies in Classical Philology, 59).}{Vérifié (EHESS ; catalogue de Berkeley) ; contenu connu par les comptes rendus.}

\bib{M.~R.~Salzman, \emph{The Making of a Christian Aristocracy. Social and Religious Change in the Western Roman Empire}, Cambridge (Mass.), Harvard University Press, 2002.}{Partiellement vérifié (BMCR 2002.06.04) ; contenu non consulté.}

\bib{W.~Scheidel, « Roman Age Structure: Evidence and Models », \emph{Journal of Roman Studies}, 91, 2001, p.~1–26.}{Partiellement vérifié ; contenu non consulté.}

\bib{B.~D.~Shaw, « Seasons of Death: Aspects of Mortality in Imperial Rome », \emph{Journal of Roman Studies}, 86, 1996, p.~100–138, DOI 10.2307/300425.}{Vérifié (Cambridge Core) ; contenu non consulté.}

\bib{K.~J.~Shelton, \emph{The Esquiline Treasure}, Londres, 1981.}{Notice vérifiée ; contenu non consulté.}

\bib{P.~Testini, \emph{Archeologia cristiana}, Rome, 1958 ; 2\textsuperscript{e} éd., Bari, 1980.}{Notice vérifiée ; contenu non consulté.}

\bib{D.~E.~Trout, \emph{Damasus of Rome. The Epigraphic Poetry. Introduction, Texts, Translations, and Commentary}, Oxford, Oxford University Press, 2015 (Oxford Early Christian Texts).}{Partiellement vérifié (pagination divergente selon les notices) ; contenu non consulté.}

\bib{V.~Väänänen, \emph{Introduction au latin vulgaire}, 3\textsuperscript{e} éd., Paris, Klincksieck, 1981.}{Notice vérifiée ; contenu non consulté.}

\section*{Comptes rendus cités}\addcontentsline{toc}{section}{Comptes rendus cités}

\bib{A.~Chastagnol, compte rendu de \emph{Consuls of the Later Roman Empire}, \emph{Journal of Roman Archaeology}, 1, 1988, p.~203–206.}{Référence transmise ; non réexaminée.}

\bib{H.~Chadwick, compte rendu de \emph{Roma christiana}, \emph{Journal of Ecclesiastical History}, 30, 1979, p.~477–479.}{Référence transmise ; non réexaminée.}

\bib{M.~Costambeys, compte rendu de \emph{Religion et sépulture}, \emph{sehepunkte}, 5, 2005, n\textsuperscript{o}~11.}{Consulté pour la révision antérieure ; non relu.}

\bib{J.~Harries, M.~Whitby, compte rendu de \emph{Consuls of the Later Roman Empire}, \emph{Classical Review}, 39, 1989.}{Titre et pagination à compléter.}

\bib{J.~R.~Martindale, compte rendu de \emph{Consuls of the Later Roman Empire}, \emph{Journal of Roman Studies}, 79, 1989.}{Pagination non donnée.}

\bib{P.~Nautin, compte rendu de \emph{Roma christiana}, \emph{Revue de l’histoire des religions}, 198, 1981.}{Pagination non donnée.}

\bib{J.~O’Donnell, compte rendu de M.~R.~Salzman, \emph{The Making of a Christian Aristocracy}, \emph{Bryn Mawr Classical Review}, 2002.06.04.}{Vu en extrait.}

\bib{Compte rendu de É.~Rebillard, \emph{Religion et sépulture}, \emph{Bryn Mawr Classical Review}, 2004.07.52.}{Nom de l’auteur non contrôlé.}

\bib{J.-M.~Sansterre, compte rendu de \emph{Roma christiana}, \emph{L’Antiquité classique}, 47, 1978.}{Pagination non donnée.}
